\partsection{B}{Cardinal obstructions to representability}

We write $\mathsf B G$ for the category with one object and endomorphism
group $G$, with categorical composition given by group multiplication, and
$BG=B(\mathsf B G)$ for its classifying space.  All groups constructed below
are $\U$-small.  Our cardinal parameters are infinite cardinals in $\U$;
there is no passage to a group outside the universe of the source category.
For every $\U$-small set $S$, an infinite cardinal $\lambda$ in $\U$ with
$\lambda>|S|$ exists: one may take the cardinal of a sufficiently large
power set of $S\amalg\N$.  Free products, group presentations on such sets,
quotients, finite products, and countable unions remain in $\U$.

\begin{claim}{B.1}{Proved}{What a sufficiently large group sees in a small category}
Let $K$ be a $\U$-small category and let $G$ be a $\U$-small group such that
every group homomorphism $G\longrightarrow H$ is trivial whenever
$|H|\leq |\Mor K|$.  The constant-functor map and evaluation map
\[
 \delta_K:K\longrightarrow\Fun(\mathsf B G,K),\qquad
 \operatorname{ev}_*: \Fun(\mathsf B G,K)\longrightarrow K
\]
are mutually inverse isomorphisms of categories.  Consequently the map
\[
 p^*:\Hom_{\Cathot}(e,K)\longrightarrow
       \Hom_{\Cathot}(\mathsf B G,K),\qquad p:\mathsf B G\longrightarrow e,
\]
is a bijection.  Both assertions are natural in $K$ among categories
whose automorphism groups receive only trivial homomorphisms from $G$.
\end{claim}
\begin{proof}
A functor $F:\mathsf B G\to K$ chooses an object $k=F(*)$ and a
homomorphism $G\to\operatorname{Aut}_K(k)$.  Invertibility is preserved by
every functor, even when $K$ has noninvertible endomorphisms or nontrivial
idempotents.  Since
$|\operatorname{Aut}_K(k)|\leq |\Mor K|$, this homomorphism is trivial.
Thus $F=\delta_K(k)$ as an actual functor.  A natural transformation
$\delta_K(k)\Rightarrow\delta_K(l)$ consists of one arrow $k\to l$ of
$K$; its naturality equations are identities.  Evaluation therefore gives
a bijection on objects and on every hom-set, and preserves composition.
If $K$ is empty, both categories in the assertion are empty and the same
conclusion holds.  Passing to components proves the last assertion.
\end{proof}

\begin{claim}{B.2}{Known}{The finite Schur-multiplier input}
For a finite set $F$ with $|F|\geq8$,
\[
 H_1(\operatorname{Alt}(F);\Z)=0,
 \qquad H_2(\operatorname{Alt}(F);\Z)\cong\Z/2.
\]
If $a,b,c,d$ are distinct elements of $F$, the commuting permutations
$x=(ab)(cd)$ and $y=(ac)(bd)$ define a torus class
\[
 t_F=[x\mid y]-[y\mid x]\in H_2(\operatorname{Alt}(F);\Z)
\]
which generates this group.  In particular, inclusions between finite
alternating groups on sets containing these four elements and having at
least eight elements carry generator to generator.
\end{claim}
\begin{proof}[Source and the required transition-map check]
The multiplier calculation is Schur's classical finite calculation.  We
use the explicitly recorded formula and generator in
\citet[Eq.~(35), Proposition~20, and Corollary~21]{BC_DominguezSegovia2021};
the cited numbering is that of arXiv version~2, rather than the later
journal version.  This is the one imported finite group calculation in
the proof of Lead~1.  We do not infer a direct-limit calculation merely
from the abstract isomorphism $H_2\cong\Z/2$: the cycle displayed above
uses the same four permutations before and after any inclusion, and its
image is the displayed nonzero generator.  This verifies the required
transition maps.

There is also a direct check of its nontriviality.  In the real Clifford
algebra on orthonormal $e_a,e_b,e_c,e_d$, put
\[
 \widetilde x=\tfrac12(e_a-e_b)(e_c-e_d),\qquad
 \widetilde y=\tfrac12(e_a-e_c)(e_b-e_d).
\]
These are spin lifts of $x$ and $y$.  They satisfy
$\widetilde x^2=\widetilde y^2=-1$ and
$\widetilde x\widetilde y=-\widetilde y\widetilde x$.
Thus the central double cover evaluates nontrivially on the torus cycle.
The exact rational Clifford-algebra verification is included in
\texttt{code/bc\_checks.py}.  This check proves nontriviality; the assertion
that the whole finite multiplier has order two is the cited classical
input.
\end{proof}

\begin{claim}{B.3}{Proved}{Finitary alternating groups at every infinite cardinal}
Let $\Omega$ be an infinite $\U$-small set and let
$A_\Omega=\operatorname{Alt}_{\mathrm{fin}}(\Omega)$ be the group of even
permutations with finite support.  Then
\[
 |A_\Omega|=|\Omega|,\qquad
 A_\Omega\text{ is nonabelian simple},\qquad
 H_1(A_\Omega;\Z)=0,\quad H_2(A_\Omega;\Z)=\Z/2.
\]
Every homomorphism from $A_\Omega$ to a group of cardinality strictly
less than $|\Omega|$ is trivial.
\end{claim}
\begin{proof}
There are $|\Omega|$ finite subsets of $\Omega$, and each supports finitely
many permutations.  This gives the upper bound on the cardinality.  The
3-cycles using two fixed elements and a varying third element give the
opposite bound.

Let $N$ be a nontrivial normal subgroup and choose $g\in N\setminus\{1\}$.
Choose $a$ moved by $g$, put $b=g(a)$, and choose $c,d$ outside the support
of $g$ and distinct from $a,b$.  With $\tau=(a\,c\,d)$, the commutator
\[
 g\tau g^{-1}\tau^{-1}=(b\,c\,d)(a\,d\,c)=(a\,b\,c)
\]
belongs to $N$.  Every 3-cycle is conjugate to this one by an even
finitary permutation: a finite permutation carrying the ordered triple
to the desired triple can have its parity adjusted on two fresh points.
Every even finitary permutation is a product of 3-cycles, by decomposing
it into an even number of transpositions and replacing pairs of
transpositions by 3-cycles.  Hence $N=A_\Omega$.  Noncommuting 3-cycles
show that the group is nonabelian.  Its commutator subgroup is therefore
a nontrivial normal subgroup, so the group is perfect and $H_1=0$.

Fix four elements $a,b,c,d$ of $\Omega$.  The finite subsets $F$ containing
these elements and with $|F|\geq8$ form a filtered cofinal subsystem of
the finite subsets of $\Omega$.  The normalized bar-chain complex of
$A_\Omega$ is the filtered colimit of those of $\operatorname{Alt}(F)$:
every finite tuple of permutations, and every finite chain, has support
in one finite $F$.  Filtered colimits of abelian groups are exact, so
\[
 H_2(A_\Omega;\Z)
   =\colim_F H_2(\operatorname{Alt}(F);\Z).
\]
By B.2 this is the constant system $\Z/2$ when its generators are chosen
to be $t_F$.  It therefore has colimit $\Z/2$.  Finally, a nontrivial
homomorphism from a simple group is injective.  Such an injection into a
strictly smaller group is impossible.
\end{proof}

\begin{claim}{B.4}{Refuted}{The asserted right adjoint and the first proposed envelope}
There is no $\U$-small category $K$ representing
\[
 X\longmapsto [BX,K(\Z/2,2)]
 \quad\text{on }\Cathot^{\op}.
\]
Consequently $\varphi:\Cathot\to\Hot$ has no right adjoint, fully
faithful or otherwise.  The assertion at V~ms.~37 and the unrestricted
Thomason--Kan-envelope construction at V~ms.~42--43 are false.
\end{claim}
\begin{proof}
Suppose there were natural bijections
$\theta_X:\pi_0\Fun(X,K)\to[BX,K(\Z/2,2)]$ for all $\U$-small $X$.
Taking $X=e$ shows that $\pi_0K$ is a singleton.  Choose an infinite
$\U$-cardinal $\lambda>|\Mor K|$, a set $\Omega$ of cardinality
$\lambda$, and $G=A_\Omega$.  By B.1 and B.3,
$\pi_0\Fun(\mathsf B G,K)=\pi_0K$ is a singleton.  On the other hand,
ordinary unpointed cohomological representability and the universal
coefficient theorem give
\[
\begin{split}
 [BG,K(\Z/2,2)]
 &=H^2(BG;\Z/2)\\
 &\cong\Hom_{\Z}(H_2(BG;\Z),\Z/2)
 \cong\Z/2.
\end{split}
\]
Here the omitted universal-coefficient term is
$\operatorname{Ext}^1_{\Z}(H_1(BG;\Z),\Z/2)=0$, because $G$ is perfect.
The two cohomology classes remain distinct as unpointed homotopy classes:
the target is simply connected.  We have a contradiction.  If a right
adjoint existed, its value at this homotopy type would be such a
representing category, so no right adjoint exists.
\end{proof}

\begin{claim}{B.5}{Proved}{A simple overgroup with a prescribed infinite cardinality}
Let $H$ be a $\U$-small group and let $\lambda$ be an infinite
$\U$-cardinal with $|H|\leq\lambda$.  There is an injective homomorphism
$H\hookrightarrow S$ into a nonabelian simple $\U$-small group $S$ with
$|S|=\lambda$.
\end{claim}
\begin{proof}
We give the construction to make its cardinal control explicit.  Start
with $H_0=H*F_\lambda$, so $H\hookrightarrow H_0$ and $|H_0|=\lambda$.
For an arbitrary group $L$ of cardinality $\lambda$, form
\[
 L'=L*F\{t_{g,h}:g\in L\setminus\{1\},\ h\in L\}.
\]
For each pair $(g,h)$ abbreviate $t=t_{g,h}$ and put
\[
 w_{g,h}=(gtg^{-1}t^{-1})(gt^2g^{-1}t^{-2}).
\]
Both $w_{g,h}$ and $w_{g,h}h$ have infinite order in the free factor
$L*\langle t\rangle$ of $L'$.  The first word is cyclically reduced of
syllable length eight.  For the second, cyclically moving its initial
$g$ to the end gives a cyclically reduced word of length eight when
$hg\ne1$.  When $hg=1$, a further cyclic reduction gives
\[
 g^{-1}t^{-1}gt^2g^{-1}t^{-1},
\]
which has syllable length six.  A cyclically reduced free-product word
of length greater than one has infinite order, since its positive powers
are reduced concatenations.

We may therefore adjoin stable letters $s_{g,h}$ identifying the two
infinite cyclic subgroups by
\[
 s_{g,h}^{-1}w_{g,h}s_{g,h}=w_{g,h}h.
\]
The simultaneous HNN embedding theorem ensures that $L'$ embeds in the
resulting group $E(L)$; its precise form, allowing an arbitrary family
of subgroup isomorphisms, is
\citet[Theorem~II]{BC_HNN1949}.  Equivalently, the HNN normal form says
that a nontrivial word in the base group remains nontrivial after these
stable letters are adjoined.  This theorem is applicable because each
specified map $w_{g,h}^r\mapsto(w_{g,h}h)^r$ is an isomorphism between
infinite cyclic subgroups of the same base group $L'$.  There are
$\lambda$ generators and relations, so $|E(L)|=\lambda$.

In $E(L)$, the element $h=w_{g,h}^{-1}s_{g,h}^{-1}w_{g,h}s_{g,h}$ lies
in the normal closure of $g$: each factor of $w_{g,h}$ already lies in
that normal closure.  Iterate the injective construction
$H_{r+1}=E(H_r)$, and set $S=\bigcup_{r<\omega}H_r$.  This union has
cardinality $\lambda$.  For any $g\ne1$ and $h$ in $S$, one stage
contains both, and the next stage puts $h$ in the normal closure of $g$.
Every nontrivial normal subgroup of $S$ is consequently $S$.  The
embedded free group $F_\lambda$ makes $S$ nonabelian.  This also proves
that none of the inclusions has collapsed the original group $H$.
\end{proof}

\begin{claim}{B.6}{Proved}{An acyclic overgroup without a cardinal increase}
For every infinite $\U$-small group $H$ there is an injective
homomorphism $H\hookrightarrow M(H)$ with
\[
 |M(H)|=|H|,\qquad \widetilde H_i(BM(H);\Z)=0
 \quad\text{for every }i\geq0.
\]
For a finite $H$, the same construction gives a countable overgroup.
\end{claim}
\begin{proof}
We use the mitosis construction of
\citet[Proposition~4.1, Theorem~4.2, and Lemma~5.4]{BC_BDH1980}.
Its algebraic form also supplies a short verification of the size bound.
Embed $H$ as the first factor of $H\times H$, and adjoin two HNN stable
letters $d,s$ satisfying
\[
 d^{-1}(h,1)d=(h,h),\qquad s^{-1}(h,1)s=(1,h)
 \quad(h\in H).
\]
Both identifications are isomorphisms between subgroups of $H\times H$.
The base group embeds, and the resulting group $m(H)$ is generated by
$H,d,s$, because the second factor is conjugate to the first.  In the
notation $h^s=s^{-1}hs$ it satisfies
\[
 h^d=hh^s,\qquad [h',h^s]=1\quad(h,h'\in H).
\]
Put $M(H)=\bigcup_{r<\omega}m^r(H)$ using these injections.  Its
cardinality is at most $\max(|H|,\aleph_0)$, and it contains $H$.

For completeness, the homology-annihilation argument is as follows.
Suppose $f:A\to B$ induces zero in field homology in positive degrees
less than $n$, and $j:B\hookrightarrow m(B)$ is the mitosis inclusion.
The homomorphism
\[
 k:B\times B\longrightarrow m(B),\qquad k(b,c)=b\,c^s
\]
is defined because the two displayed subgroups commute.  The two factor
inclusions followed by $k$ are conjugate to $j$, and the diagonal
followed by $k$ is also conjugate to $j$.  For $z\in H_n(A;k)$ the
K\"unneth formula and the vanishing of the lower-degree maps give
\[
 (f\times f)_*\Delta_*(z)=f_*(z)\otimes1+1\otimes f_*(z).
\]
Applying $k_*$ and using invariance under inner automorphisms yields
$(jf)_*(z)=2(jf)_*(z)$, hence $(jf)_*(z)=0$ over every field.  Induction
shows that a composite of $n$ mitosis inclusions kills field homology
through degree $n$.  Every bar cycle in $M(H)$ occurs at a finite
stage, so all its positive field homology vanishes.  Rational vanishing
makes each positive integral homology group torsion.  The universal
coefficient sequence in the next degree, with coefficients $\F_p$,
shows that it has no $p$-torsion for any prime $p$.  It is therefore zero.
The group is nonempty, so its classifying space is connected as required.
\end{proof}

\begin{claim}{B.7}{Proved}{Acyclic simple overgroups of arbitrary size}
For every group $H$ and every infinite $\U$-cardinal
$\lambda\geq\max(|H|,\aleph_0)$ there is a monomorphism of groups
$i:H\hookrightarrow A$ such that $A$ is nonabelian simple,
$|A|=\lambda$, and $BA$ is integrally acyclic.
\end{claim}
\begin{proof}
Start with $H_0=H*F_\lambda$.  Alternate B.5 and B.6 to obtain
injective homomorphisms
\[
 H_0\hookrightarrow S_0\hookrightarrow M_0\hookrightarrow
 S_1\hookrightarrow M_1\hookrightarrow\cdots,
\]
where every $S_r$ is simple, every $M_r$ is acyclic, and every term has
cardinality $\lambda$.  Let $A$ be the union.  If a normal subgroup
$N\lhd A$ contains $g\ne1$, then for any $a\in A$ a sufficiently late
$S_r$ contains both $g$ and $a$.  The subgroup $N\cap S_r$ is normal in
$S_r$ and nontrivial, hence contains $a$.  This proves simplicity.
Every finite bar cycle lies in some $M_r$ and bounds there, so $A$ is
acyclic.  Its cardinality is $\lambda$, and the inclusions give the
asserted monomorphism from $H$.  The alternating-union observation is a
formal combination of the two preceding constructions; no priority is
claimed for it.
\end{proof}

\begin{claim}{B.8}{Proved}{Groups invisible to smaller groups and mapping to spheres}
For every integer $n\geq2$ and every infinite $\U$-cardinal $\lambda$
there is a $\U$-small group $G_{n,\lambda}$ and a pointed continuous map
of CW models
\[
 q_{n,\lambda}:BG_{n,\lambda}\longrightarrow S^n
\]
such that $|G_{n,\lambda}|=\lambda$, the map $q_{n,\lambda}$ induces
isomorphisms on integral homology in all degrees, and every homomorphism
from $G_{n,\lambda}$ to a group of cardinality less than $\lambda$ is
trivial.  The group $G_{n,\lambda}$ is perfect.  It is not asserted to be
simple.
\end{claim}
\begin{proof}
For a group $H$ with $|H|\leq\lambda$, choose the embedding
$i:H\hookrightarrow A$ of B.7 and form the amalgamated free product
\[
 D_\lambda(H)=A*_{H}A.
\]
The maps from both factors $A$ to this amalgam are injective by the
amalgamated-product normal form.  Thus $|D_\lambda(H)|=\lambda$.
A homomorphism $D_\lambda(H)\to L$ with $|L|<\lambda$ restricts
trivially to each factor, because each factor is simple of cardinality
$\lambda$.  Since the two factors generate the amalgam, the homomorphism
itself is trivial.

The double mapping cylinder of
$BA\leftarrow BH\rightarrow BA$, whose structure maps are $Bi$, is a
classifying CW space for $D_\lambda(H)$.  To see the asphericity required
here, the universal cover is the tree of spaces over the Bass--Serre
tree of the amalgam, with contractible vertex and edge spaces.  A finite
subtree gives a contractible union, by induction over its terminal
vertices; every sphere and homotopy meets a finite subtree after cellular
approximation.  The total cover is therefore weakly contractible.  The
ordinary van Kampen theorem identifies the fundamental group with the
specified amalgam.  This establishes the comparison with a classifying
space without substituting a homotopy pushout for an ordinary pushout
without justification.

Collapse the two copies of $BA$ in this double mapping cylinder to the
two suspension vertices.  We obtain a pointed map
\[
 r_H:B D_\lambda(H)\longrightarrow \Sigma BH.
\]
Here $\Sigma$ denotes the unreduced suspension of the connected CW
space $BH$, with one suspension vertex chosen as basepoint; it agrees
up to pointed homotopy with its reduced suspension after choosing the
standard vertex of $BH$.  The Mayer--Vietoris sequence, compared with
that for the suspension, proves that $r_H$ is an integral homology
isomorphism, because both copies of $BA$ are acyclic.  In particular
$H_1(D_\lambda(H);\Z)=0$.

Start with $G_{1,\lambda}=\Z$ and a chosen pointed homotopy equivalence
$q_{1,\lambda}:B\Z\to S^1$.  Inductively set
\[
 G_{n,\lambda}=D_\lambda(G_{n-1,\lambda}),\qquad
 q_{n,\lambda}=(\Sigma q_{n-1,\lambda})r_{G_{n-1,\lambda}},
\]
using the usual pointed identification $\Sigma S^{n-1}\simeq S^n$.
Suspension preserves integral homology isomorphisms, so induction proves
the homology assertion.  Every group with $n\geq2$ is the specified
double and therefore has all the remaining properties already checked.
\end{proof}

\begin{claim}{B.9}{Proved}{The target condition needed to use these homology equivalences}
Let $u:A\to B$ be a pointed integral homology equivalence between
connected CW spaces, and let $Z$ be a pointed simply connected CW space.
Precomposition induces a weak equivalence of derived pointed mapping
spaces
\[
 u^*:\Map_*(B,Z)\longrightarrow\Map_*(A,Z).
\]
In particular it induces a bijection $[B,Z]_*\to[A,Z]_*$.
\end{claim}
\begin{proof}
For an abelian group $M$, the universal coefficient theorem gives an
isomorphism $\widetilde H^j(B;M)\to\widetilde H^j(A;M)$ for every
$j\geq0$.  The homotopy groups of
$\Map_*(A,K(M,r))$, in every component, are
$\widetilde H^{r-i}(A;M)$ for $0\leq i\leq r$, and vanish for $i>r$;
the mapping space is an abelian group object.  Hence the assertion holds
for every $K(M,r)$.

For a simply connected $Z$, choose its Postnikov tower with
$P_1Z=*$ and with $P_rZ$ the homotopy fibre of its $k$-invariant
\[
 P_{r-1}Z\longrightarrow K(\pi_rZ,r+1).
\]
All coefficient systems in this assertion are constant, precisely
because $Z$ is simply connected.  Derived mapping spaces preserve
homotopy limits in the target.  Induction on $r$, using the homotopy
fibre descriptions, proves the assertion for every $P_rZ$.  Finally
$Z\to\operatorname{holim}_rP_rZ$ is a weak equivalence: for each
homotopy group the tower is eventually constant with the value of that
homotopy group of $Z$.  Taking the homotopy limit of the already proved
mapping-space equivalences gives the assertion for $Z$.  This last step
is essential; we have proved equivalences of mapping spaces at all
stages, rather than merely bijections of their component sets, so no
uncontrolled inverse-limit or phantom-map inference occurs.  The result
is the integral-homology locality of simply connected spaces; see
\citet[Example~9.2]{BC_Lawson2019}.
\end{proof}

\begin{claim}{B.10}{Proved}{Essential maps from arbitrarily large groups for every higher type}
Let $Y$ be a $\U$-small homotopy type with a point $y$ and a nonzero
element $\alpha\in\pi_n(Y,y)$ for some $n\geq2$.  For every infinite
$\U$-cardinal $\lambda$, the group $G_{n,\lambda}$ of B.8 has no
nontrivial homomorphism to any group of cardinality less than $\lambda$
and admits a map $BG_{n,\lambda}\to Y$ which is not freely homotopic to
a constant map.  Thus this assertion holds for every homotopy type that
is not a $1$-type, including every simply connected noncontractible
homotopy type.
\end{claim}
\begin{proof}
Choose a CW model for the component $Y_y$ and its pointed universal
cover $p:\widetilde Y_y\to Y_y$.  The sphere map representing $\alpha$
lifts uniquely as a pointed map
$\widetilde\alpha:S^n\to\widetilde Y_y$, and its class is nonzero
because a covering induces isomorphisms in homotopy degrees at least
two.  By B.9 and B.8,
$\widetilde\alpha q_{n,\lambda}:BG_{n,\lambda}\to\widetilde Y_y$
is not pointed nullhomotopic.

The composite $p\widetilde\alpha q_{n,\lambda}$ cannot be freely
nullhomotopic in $Y$.  If a free nullhomotopy existed, its image would
lie in $Y_y$ and its lift beginning at
$\widetilde\alpha q_{n,\lambda}$ would end at a constant map in the
cover.  The lifted homotopy exists by the covering homotopy property.
Its terminal map is constant because its image lies in the discrete
fibre over the endpoint and $BG_{n,\lambda}$ is connected.  In a
simply connected target, a free nullhomotopy of a pointed map can be
made pointed: the obstruction is the action of the target fundamental
group, which is trivial.  We would obtain a pointed nullhomotopy,
contrary to B.9.  This proves essentiality in the unpointed sense
actually required.  The last sentence follows by choosing a nonzero
higher homotopy group; a connected simply connected CW type with all
homotopy groups zero is contractible by Whitehead's theorem.
\end{proof}

\begin{claim}{B.11}{Proved}{The full negative representability statement}
For every $\U$-small homotopy type $Y$ which is not a $1$-type, the
presheaf $X\mapsto[BX,Y]$ on $\Cathot$ is not representable by a
$\U$-small category.  No connectedness assumption on $Y$ is needed.
\end{claim}
\begin{proof}
Suppose that $K$ represents the presheaf, by natural bijections
$\theta_X:\pi_0\Fun(X,K)\to[BX,Y]$.  Choose
$\lambda>|\Mor K|$ in $\U$, and choose $n$ and the corresponding
$G=G_{n,\lambda}$ from B.10.  Naturality for the unique functor
$p:\mathsf B G\to e$ gives a commutative square
\[
\begin{array}{ccc}
\pi_0K&\xrightarrow{\ p^*\ }&\pi_0\Fun(\mathsf B G,K)\\
\theta_e\downarrow\phantom{\theta_e}&&\phantom{\theta_{\mathsf B G}}\downarrow\theta_{\mathsf B G}\\
\pi_0Y&\xrightarrow{\ (Bp)^*\ }&[BG,Y].
\end{array}
\]
The upper arrow is a bijection by B.1.  The lower arrow sends a component
of $Y$ to its constant-map class.  It is not surjective by B.10.  This
contradicts the vertical bijections.  This argument uses the specified
constant-map function, rather than the cardinalities of the two sets;
it therefore also applies when $Y$ has infinitely many components.
\end{proof}

\partsection{C}{Local categories and the groupoid reflection}

For a category $X$, let $\Pi_1X=X[\Mor(X)^{-1}]$ denote the groupoid
obtained by inverting all its arrows, and let
$\eta_X:X\to\Pi_1X$ be the canonical functor.  Its objects are those
of $X$; its arrows are finite composable words of arrows and formal
inverses, modulo the identity, inverse, and composition relations.
This construction is $\U$-small.  It is equivalent to the fundamental
groupoid of $BX$ on the vertices of its nerve: the nerve's edges give
the generating paths, its $2$-simplices impose the composition
relations, and higher cells impose no further relations on the
fundamental groupoid.  We allow the empty groupoid.

\begin{claim}{C.1}{Proved}{Groupoids are local, with the canonical comparison}
For every $\U$-small category $X$ and $\U$-small groupoid $G$, the
canonical map
\[
 \pi_0\Fun(X,G)\longrightarrow[BX,BG]
\]
is a bijection, natural in both variables.  More precisely,
\[
 N\Fun(X,G)\cong\operatorname{Map}_{\mathrm{sSet}}(NX,NG)
\]
is an isomorphism of simplicial sets, and the right side is a derived
mapping space.  Restriction along $\eta_X$ is an isomorphism of
categories $\Fun(\Pi_1X,G)\to\Fun(X,G)$.
\end{claim}
\begin{proof}
In degree $r$, the asserted simplicial isomorphism is the chain of
canonical bijections
\[
\begin{split}
N\Fun(X,G)_r
 &=\Fun([r],\Fun(X,G))\\
 &=\Fun(X\times[r],G)\\
 &=\Hom_{\mathrm{sSet}}(NX\times\Delta[r],NG).
\end{split}
\]
The last equality uses full faithfulness of the nerve and preservation
of products.  These bijections commute with faces and degeneracies.
The nerve of a groupoid is a Kan complex, and every simplicial set is
cofibrant in the Kan--Quillen model structure.  Thus the displayed
mapping space already computes derived maps of spaces, giving precisely
the canonical map on components.  This proves localness, including for
$X$ or $G$ empty.

Every functor $X\to G$ extends uniquely over $\eta_X$, since its values
on arrows are invertible.  A natural transformation between two such
functors has invertible components.  Its naturality square for an arrow
of $X$ also gives the naturality square for its formal inverse, by
multiplying the square by the inverses.  Therefore transformations extend
uniquely as well, and restriction is the asserted isomorphism of
categories.  This last assertion concerns ordinary categories, before
passing to components.
\end{proof}

\begin{claim}{C.2}{Proved}{Classification of all representable homotopy types}
A $\U$-small homotopy type $Y$ has its presheaf
$X\mapsto[BX,Y]$ represented on $\Cathot$ by a $\U$-small category if
and only if $Y$ is a $1$-type.  If $Y$ is a $1$-type, a fundamental
groupoid $G_Y$ represents it.  Every other representing category is
isomorphic to $G_Y$ in $\Cathot$, and its classifying space is
isomorphic to $Y$ in $\Hot$.  This is a complete affirmative answer to
Lead~2's representability conjecture.
\end{claim}
\begin{proof}
The negative direction is B.11.  If $Y$ is a $1$-type, choose its
fundamental groupoid, with one representative object for each component
and automorphism group $\pi_1(Y,y)$ there.  This is $\U$-small, and
$BG_Y\simeq Y$: the comparison identifies components and fundamental
groups, and both sides have zero higher homotopy groups.  C.1 gives
the required natural bijections.

If a category $K$ represents the same presheaf, the resulting natural
isomorphism
$\Hom_{\Cathot}(-,K)\cong\Hom_{\Cathot}(-,G_Y)$ is represented by an
isomorphism $K\cong G_Y$ in $\Cathot$, by the ordinary Yoneda lemma.
Applying $B$ gives an isomorphism in $\Hot$.  More precisely, if
$\theta$ is the specified representing isomorphism, its universal
element $\theta_K(\id_K):BK\to Y$ is this weak equivalence in
$\Hot$, composed with the chosen identification $BG_Y\simeq Y$.
Thus an unspecified abstract bijection has not been substituted for the
canonical comparison.  When $Y$ is empty the representing category is
the empty groupoid, and the same Yoneda argument applies.
\end{proof}

\begin{claim}{C.3}{Proved}{Classification of local categories}
For a $\U$-small category $K$, the following conditions are equivalent:
the canonical function $\pi_0\Fun(X,K)\to[BX,BK]$ is bijective for
every $\U$-small $X$; the category $K$ is isomorphic in $\Cathot$ to a
groupoid; and the canonical functor $\eta_K:K\to\Pi_1K$ is an
isomorphism in $\Cathot$.  In particular every local $K$ has $BK$ a
$1$-type.  Conversely, for a category with $BK$ a $1$-type, localness
is equivalent to being isomorphic in $\Cathot$ to its groupoidification.
The hypothesis that $BK$ is a $1$-type alone is insufficient.
\end{claim}
\begin{proof}
If $K$ is local, it represents $X\mapsto[BX,BK]$.  C.2 implies that
$BK$ is a $1$-type.  Hence $B\eta_K:BK\to B\Pi_1K$ is a weak
equivalence: it always induces isomorphisms on components and
fundamental groups, and in this case both source and target have no
higher homotopy groups.

Localness at the source $X=\Pi_1K$ produces a morphism
$[j]:\Pi_1K\to K$ in $\Cathot$ whose realization represents
$(B\eta_K)^{-1}$ in $\Hot$.  Localness of $K$ at $X=K$ then implies
$[j][\eta_K]=\id_K$.  C.1, applied to the groupoid $\Pi_1K$, gives
$[\eta_K][j]=\id_{\Pi_1K}$.  These are equalities in $\Cathot$;
they do not assert actual equalities of functors or an equivalence of
the underlying categories.

The third condition implies the second.  To prove the converse, we
observe that postcomposition by an isomorphism $K\to G$ in $\Cathot$
induces bijections on all representable hom-sets, and realization of
that isomorphism induces bijections $[BX,BK]\to[BX,BG]$.  The square
formed by the two canonical comparisons commutes, and the comparison
for $G$ is bijective by C.1.  Thus the comparison for $K$ is bijective.
The failure of the weaker converse is exhibited in C.5 below.
\end{proof}

\begin{claim}{C.4}{Proved}{The reflector is fundamental groupoidification}
Let $\mathcal L\subset\Cathot$ be the full subcategory of local
objects.  It is reflective.  After identifying it with groupoids modulo
natural isomorphism, its reflector is
\[
 \Pi_1:\Cathot\longrightarrow\operatorname{Ho}(\mathrm{Grpd}),
 \qquad \Pi_1\dashv I,
\]
where $I$ is the full inclusion and the unit is $[\eta_X]$.
This reflector inverts exactly those functors that induce a bijection
on components and isomorphisms on fundamental groups at every
basepoint.  Its unit is in $W_\infty$ precisely when $BX$ is a
$1$-type.  In particular this reflection does not supply the nonexistent
right adjoint to $\varphi:\Cathot\to\Hot$.
\end{claim}
\begin{proof}
Applying $\Pi_1$ to a natural transformation produces a natural
isomorphism.  It therefore descends to $\Cathot$.  C.1 gives natural
bijections
\[
 \Hom_{\operatorname{Ho}(\mathrm{Grpd})}(\Pi_1X,G)
   =\pi_0\Fun(\Pi_1X,G)
   \cong\pi_0\Fun(X,G)
   =\Hom_{\Cathot}(X,IG).
\]
This is the asserted adjunction, with the displayed unit.  C.3
identifies its replete essential image with $\mathcal L$.
A functor of groupoids is an equivalence exactly when it gives the
specified bijection of components and isomorphisms of isotropy groups;
thus the maps inverted are exactly the stated $1$-equivalences.
For every $X$, the map $B\eta_X$ preserves components and fundamental
groups and has a $1$-type as its target.  It is a weak equivalence if
and only if all the higher homotopy groups of $BX$ vanish.

Although every $W_\infty$-map is inverted by this reflector, many more
maps are inverted.  For example a category with classifying space $S^2$
maps to a contractible groupoidification.  The target $\mathcal L$
therefore models $1$-truncation; it does not model the localization
$\Hot$ at $W_\infty$.  This distinction is necessary when using the
word ``reflective'' here.
\end{proof}

\begin{claim}{C.5}{Refuted}{Coreflectivity, and the weaker converse from the classifying space}
The inclusion $I:\mathcal L\hookrightarrow\Cathot$ has no right
adjoint.  Moreover, a category with contractible classifying space need
not be local.
\end{claim}
\begin{proof}
Let $F$ be the poset with object set $\N$, identity arrows, and exactly
the following nonidentity arrows:
\[
 2r\longrightarrow2r+1\quad(r\geq0),\qquad
 2r\longrightarrow2r-1\quad(r\geq1).
\]
There are no composable pairs of nonidentity arrows, so composition
with identities specifies the entire category.  Its nerve has only
vertices and these nondegenerate edges.  Its realization is the
infinite ray, hence is contractible.

Let $D$ be the discrete category with object set $\N$.  Define functors
$f,g:D\to F$ by $f(r)=0$ and $g(r)=2r$.  A natural transformation
between two functors $D\to F$ is an objectwise family of arrows.  Thus
each one-step transformation, in either direction, changes each
coordinate by at most one in the undirected path metric on $F$.
A zigzag of length $m$ changes every coordinate by at most $m$.
Since the distance from $f(r)$ to $g(r)$ is $2r$, no finite zigzag
connects $f$ to $g$.  Hence
\[
 |\Hom_{\Cathot}(D,F)|>1,
 \qquad |\Hom_{\Cathot}(e,F)|=1.
\]
In particular the canonical comparison to $[BD,BF]$, which is a
singleton, is not injective, proving that $F$ is not local.  The
finite truncations and their exact distances are checked by
\texttt{code/bc\_checks.py}; the unbounded-coordinate argument proves
the infinite assertion.

Within $\mathcal L$, the object $D$ is the coproduct of countably many
copies of $e$.  Indeed groupoids have this coproduct, and a family of
componentwise natural isomorphisms can be assembled into a single
natural isomorphism on a discrete union.  If $I$ had a right adjoint,
then, being a left adjoint, it would preserve this coproduct.
Equivalently its adjunction would force the canonical restriction map
\[
 \Hom_{\Cathot}(D,F)\longrightarrow
     \prod_{r\in\N}\Hom_{\Cathot}(e,F)
\]
to be a bijection.  Its target is a singleton and its source has the
two distinct elements just constructed.  This contradiction proves
the failure of coreflectivity.  The use of an infinite source is part
of the counterexample, not a suppressed finiteness assumption.
\end{proof}

\begin{claim}{C.6}{Refuted}{Thomason fibrancy is not the manuscript's localness}
It is false that every Thomason-fibrant category is local in the sense
of C.3.  There is no change of target by a $W_\infty$-equivalence that
makes a homotopy type with nonzero higher homotopy groups into a local
category.
\end{claim}
\begin{proof}
Let $P$ be the face poset of the boundary of a tetrahedron, and choose
a Thomason fibrant replacement $P\to RP$.  By the definition of
fibrant replacement this is a specified $W_\infty$-map, and
$B(RP)\simeq BP\simeq S^2$.  If $RP$ were local, C.3 would make this
a $1$-type, a contradiction.  The same reasoning applies to any
category whose nerve has a nonzero higher homotopy group and to any
$W_\infty$-equivalent target.  Fibrancy can be used for derived
mapping objects with the appropriate source and cylinder; it cannot
produce the envelope asserted at V~ms.~37--43.
\end{proof}
